% ---------------------------------------------------------------- title page
\begin{titlepage}
\pagecolor{bg}\color{fg}
\vspace*{0.6in}
{\color{muted}\large A deep-dive tutorial on the \texttt{score} repository}\par
\vspace{10pt}
{\fontsize{38}{44}\selectfont\bfseries\color{acc} Sockets, Columns\\[2pt] and Tail Risk}\par
\vspace{14pt}
{\Large\color{fg} Designing a networked stock-risk engine from first principles:\\ TCP framing, reactors, columnar storage, Value at Risk,\\ and the floating-point bug that hid in a confidence level}\par
\vspace{26pt}
{\color{codeframe}\rule{\linewidth}{0.8pt}}\par
\vspace{12pt}
\begin{tikzpicture}
  % a small decorative frame header: 4 length bytes then payload
  \foreach \i/\b in {0/00,1/00,2/00,3/28}{\node[byte,draw=yel,text=yel] at (\i*0.86,0) {\b};}
  \foreach \i/\b in {4/7B,5/22,6/63,7/6F,8/6E,9/66,10/22,11/3A,12/30,13/2E,14/39,15/35}{\node[byte,draw=acc] at (\i*0.86,0) {\b};}
  \node[byte,draw=muted,text=muted] at (16*0.86,0) {\dots};
  \draw[decorate,decoration={brace,amplitude=5pt,mirror},draw=yel] (-0.4,-0.45) -- (3*0.86+0.4,-0.45) node[midway,below=6pt,font=\footnotesize,text=yel]{length = 40 (big-endian)};
  \draw[decorate,decoration={brace,amplitude=5pt,mirror},draw=acc] (4*0.86-0.4,-0.45) -- (15*0.86+0.4,-0.45) node[midway,below=6pt,font=\footnotesize,text=acc]{\texttt{\{"conf":0.95,...} (JSON payload)};
\end{tikzpicture}
\vspace{22pt}

\begin{minipage}[t]{0.55\linewidth}
\color{muted}\small
\textbf{\color{fg}What this is.} A textbook-style companion to the \texttt{score} C++23 library's
storage, networking and risk extensions. Every number, byte and bit pattern in it was produced by
running code; every repository listing is quoted verbatim from the working tree.\par
\vspace{6pt}
\textbf{\color{fg}How to read it.} Work each chapter in order: objectives, exposition, worked examples,
then the engineering questions. Hints, not solutions, are in \cref{app:hints}.
\end{minipage}\hfill
\begin{minipage}[t]{0.38\linewidth}
\color{muted}\small
\textbf{\color{fg}Repository state.} Branch \texttt{formatter}; newest changelog entry dated 2026-10-17.\par
\vspace{6pt}
\textbf{\color{fg}Measurements.} Benchmarks ran on an Apple M3 Max (14 cores), macOS 26.6.2, Apple clang 21.\par
\vspace{6pt}
\textbf{\color{fg}Compiled.} \today.
\end{minipage}
\vfill
{\color{muted}\footnotesize Prepared with AI assistance (Claude). The text and figures are largely AI-generated and were checked by
re-running the code that produced every quoted number; please review before relying on it.}
\end{titlepage}

% ---------------------------------------------------------------- abstract
\chapter*{Abstract}
\addcontentsline{toc}{chapter}{Abstract}
\markboth{Abstract}{}
The \texttt{score} repository began as a C++23 library of descriptive statistics and stock metrics
(log-returns, the Sharpe ratio, maximum drawdown) exposed to Python through nanobind. It had three
conspicuous gaps: nothing was ever \emph{stored} (prices arrived either from live \texttt{yfinance}
calls or from a synthetic geometric-Brownian-motion generator, and both were discarded after use);
nothing measured \emph{tail} risk (no Value at Risk, no Expected Shortfall); and there was no
\emph{networking or concurrency} code at all. A structured design interview, a ``grilling session'',
turned the vague wish ``add a database and some networking'' into twelve falsifiable decisions and a
twelve-phase roadmap.

This tutorial teaches the ideas behind those decisions from first principles. We treat TCP as what
it is, an ordered byte stream with no message boundaries, and prove that a length-prefixed decoder
reconstructs messages under every possible segmentation. We compare thread-per-connection,
\texttt{poll}, \texttt{epoll}/\texttt{kqueue} and \texttt{io\_uring} with a cost model and with
measurements taken on the author's machine. We derive the byte-level advantage of columnar storage
for the one query a risk engine runs constantly. We define Value at Risk and Expected Shortfall
precisely, relate the repository's estimator to order statistics and to the Hyndman--Fan quantile
taxonomy, and verify a counterexample to the subadditivity of VaR. Finally, we dissect a real
IEEE-754 defect found while implementing VaR, in which \(1.0-0.90\) evaluates to
\(0.09999999999999998\), bound the conditions under which the repository's \(10^{-9}\) guard is
provably correct, and set it beside three historical numerical failures.

Each chapter ends with engineering questions specific to this codebase. Citations point to primary
sources: original papers, standards, and official reports.

\begin{remarkbox}[title={Colour conventions used in every figure}]
\textcolor{acc}{Blue}: interfaces, structure, the ``current'' item.
\textcolor{grn}{Green}: implemented, correct, done.
\textcolor{yel}{Yellow}: length prefixes, headers, ``next''.
\textcolor{red}{Red}: failure, loss, the thing that breaks.
\textcolor{pur}{Purple}: Python producers and remarks.
\textcolor{org}{Orange}: worked examples, inline code, historical notes.
\end{remarkbox}

% ---------------------------------------------------------------- contents
{\hypersetup{linkcolor=fg}\setlength{\parskip}{0pt}\tableofcontents}

% ---------------------------------------------------------------- notation
\chapter*{Notation and glossary}
\addcontentsline{toc}{chapter}{Notation and glossary}
\markboth{Notation and glossary}{}
\begin{center}
\small
\begin{tabularx}{\linewidth}{@{}>{\leavevmode\color{yel}}l X@{}}
\toprule
\textcolor{muted}{Symbol} & \textcolor{muted}{Meaning}\\
\midrule
$P_t$ & closing price on trading day $t$\\
$r_t=\ln(P_t/P_{t-1})$ & daily log-return (\cc{risk::log_returns})\\
$X$, $L=-X$ & return (profit-and-loss fraction) and the corresponding loss\\
$F_X(x)=\Prob(X\le x)$ & cumulative distribution function (CDF) of $X$\\
$F_n$ & empirical CDF of a sample of size $n$: $F_n(x)=\frac1n\#\{i:X_i\le x\}$\\
$\ord{1}\le\dots\le\ord{n}$ & order statistics (the sample sorted ascending, as \cc{std::sort} leaves it)\\
$c$, $\alpha$ & confidence level, e.g.\ $0.95$ (the code's \cc{confidence_level})\\
$p=1-c$ & tail probability\\
$k=\lfloor pn+\varepsilon\rfloor$ & the code's \cc{tail_index}, $\varepsilon=10^{-9}$\\
$\VaR_c$, $\ES_c$ & Value at Risk and Expected Shortfall (CVaR) at confidence $c$\\
$\fl(x)$ & the IEEE-754 binary64 number nearest to the real $x$ (round-to-nearest-even)\\
$u=2^{-53}$ & unit roundoff of binary64\\
$\ulp(x)$ & unit in the last place: spacing of binary64 numbers near $x$\\
$n$ & number of connections (\cref{ch:conc}) or of observations (\cref{ch:risk,ch:float}); context decides\\
$\Theta$, $O$ & asymptotic tight bound and upper bound\\
\bottomrule
\end{tabularx}
\end{center}

\begin{center}
\small
\begin{tabularx}{\linewidth}{@{}>{\leavevmode\bfseries\color{acc}}p{3.2cm} X@{}}
\toprule
\textcolor{muted}{Term} & \textcolor{muted}{Definition}\\
\midrule
Byte stream & An ordered, reliable sequence of bytes with no record boundaries (TCP's service model).\\
Frame & One application message plus the framing bytes that delimit it on the wire.\\
Codec & Encoder/decoder pair turning frames into bytes and back (\cc{FrameCodec}).\\
Reactor & Event-demultiplexing pattern: the loop learns a handle is \emph{ready}, then calls the handler, which performs the I/O.\\
Proactor & The loop \emph{starts} I/O on the handler's behalf and calls it on \emph{completion}.\\
Readiness API & \cc{select}, \cc{poll}, \cc{epoll}, \cc{kqueue}: report which descriptors can make progress.\\
Completion API & \cc{io_uring} (Linux), IOCP (Windows): report which submitted operations finished.\\
Column store & Storage that keeps each attribute's values contiguous rather than each record's.\\
OHLCV & Open, high, low, close, volume: one daily bar (\cc{score::db::Bar}).\\
Historical simulation & Estimating risk from the empirical distribution of past returns, without a parametric model.\\
Coherent risk measure & A risk measure satisfying the four axioms of Artzner et al.\ (\cref{def:coherent}).\\
Scaffolding & An interface plus stub bodies that throw \cc{score::NotImplementedError}.\\
\bottomrule
\end{tabularx}
\end{center}


